Thresholding is one of the most straightforward and widely used approaches for detecting blink activity from EEG or EOG signals~\citep{tran2021detection,chang2016detection}.
The simplicity of thresholding makes it computationally efficient and easy to implement, particularly in real-time fatigue-detection systems~\citep{tran2021detection,jefri2022eye}. In addition, thresholding is often used as an initial detection step in more complex pipelines, where the detected candidates are further refined using morphological rules or machine-learning classifiers~\citep{jiang2023eyeblink,ghosh2023automatic}. However, selecting an appropriate threshold value is not trivial~\citep{guttmann2019new,wang2022multidimensional}. Blink amplitude can vary across participants, electrode locations, recording devices, signal quality, and experimental conditions~\citep{guttmann2019new,wang2022multidimensional}. A threshold that is too low may produce false detections due to noise or other artifacts, whereas a threshold that is too high may miss genuine blink events~\citep{tran2021detection,jefri2022eye}. Therefore, threshold selection remains a critical consideration in the reliable detection of blink activity for fatigue-driving research~\citep{ko2020eyeblink,shahbakhti2023fusion}.



% with later stages rather than treating it as the final detector; however, this makes the thresholding stage itself worthy of focused study because its quality propagates downstream.

%Thresholding remains a foundational stage in eye-blink detection because it is widely used both as a first line of defence and as a candidate-generation or preprocessing step that precedes more complex machine-learning-based and morphology-based pipelines \citep{chang2016detection,kleifges2017blinker,tran2021detection,zhang2023method,agarwal2019blink}. Some pipelines use thresholding as a candidate-generation step that precedes later refinement, refining its candidates
